\chapter{Gastric Balloon Insertion}

\section*{Overview}
Gastric balloon insertion places a temporary, soft silicone balloon inside the stomach to help with weight loss. The balloon is positioned through the mouth using an endoscope and is then filled with saline or air. It remains in place for 6 to 12 months and is later removed endoscopically.

\section*{Alternative Names}
Intragastric balloon, gastric balloon therapy, endoscopic balloon weight loss device

\section*{Principle}
The inflated balloon takes up space in the stomach, reducing its capacity for food and increasing the sensation of fullness after small meals. This early satiety helps the person eat less and adopt more sustainable eating habits. Because the balloon is temporary, it is intended to support lifestyle change rather than to act as a permanent solution.

\section*{History}
The concept of an intragastric balloon dates to the early 1980s, with modern inflatable devices approved in the 2010s. Balloon designs and placement techniques have evolved to improve safety, tolerability, and weight loss results. Today it is a recognized nonsurgical option in the management of obesity.

\section*{Why This Test or Procedure Is Ordered}
It is offered to adults with a body mass index (BMI) in the moderate obesity range who have not achieved lasting weight loss through diet and exercise alone. It may be used when the person does not qualify for or does not desire bariatric surgery. Weight loss from the balloon can also improve conditions linked to obesity, such as type 2 diabetes and high blood pressure.

\section*{Is This a Primary or Secondary Test or Procedure}
It is a secondary procedure, used when lifestyle measures alone have been insufficient. It may also serve as a bridge to more definitive bariatric surgery by helping people lose weight and reduce surgical risk beforehand.

\section*{How This Test or Procedure Is Performed}
Under sedation, an endoscope is passed through the mouth and esophagus into the stomach, and the deflated balloon is placed and then inflated with saline or air. The procedure typically takes 20 to 30 minutes and usually does not require an overnight stay. The balloon is removed with the endoscope after several months.

\section*{What Preparation Is Needed}
You must fast for several hours before the procedure. Medications that increase bleeding risk may need to be paused, and your doctor will review any prior stomach surgery or conditions. A blood test and discussion of your medical history are usually completed in advance.

\section*{Contraindications and Precautions}
\begin{itemize}
  \item Prior stomach or esophageal surgery, active peptic ulcer disease, or gastrointestinal bleeding.
  \item Pregnancy or breastfeeding.
  \item Severe reflux disease, large hiatal hernia, or an inflamed esophagus that has not been evaluated.
\end{itemize}

\section*{Risks}
\begin{itemize}
  \item Nausea, vomiting, and abdominal cramping, which are common in the first days after placement.
  \item Stomach ulceration, irritation, or bleeding from the balloon surface.
  \item Balloon deflation, which risks blockage of the intestine (obstruction) and requires prompt removal.
\end{itemize}

\section*{What Happens After the Test or Procedure}
Nausea and discomfort are expected for the first few days and are managed with medication. A structured diet program and exercise plan are essential to maximize and maintain weight loss. The balloon is removed at the planned time, usually 6 to 12 months after insertion.

\section*{Understanding Your Results}
Weight loss is gradual and varies with diet, activity, and adherence to the program. Many people lose a meaningful portion of their excess weight over the life of the balloon. Results depend on sustaining healthy habits after removal, since the balloon alone is temporary.

\section*{Normal Results}
A successful outcome is safe placement and removal of the balloon with significant, sustained weight loss. Improvement in obesity-related conditions, such as blood sugar and blood pressure, is also a marker of success.

\section*{Follow-up Tests and Procedures}
Regular visits with the weight management team monitor weight, nutrition, and tolerance of the balloon. Bariatric surgery may be discussed after removal if further weight loss is needed. Routine follow-up continues for at least a year to support long-term weight maintenance.

\section*{References}
\begin{enumerate}
  \item MedlinePlus. Gastric Balloon. U.S. National Library of Medicine; 2025.
  \item Current Medical Diagnosis and Treatment. McGraw-Hill; 2025.
\end{enumerate}

\clearpage
